\section{Benchmarking methodology for data systems}
\label{related-work:benchmarking-methodology-for-data-systems}

\textcite{beyer_2015_benchmarking_and_resource_measurement} identifies a set of requirements for reliable benchmarking, including accurate measurement of CPU time and memory, reliable termination of processes, and a controlled execution environment. They show that classical Linux mechanisms such as \texttt{time} and \texttt{ulimit} fail to meet these requirements, primarily because they do not account for child processes correctly, and propose a framework based on Linux control groups (cgroups), a kernel feature that tracks and limits resource usage for groups of processes as a unit. The \acrfull{jcgm} supplies the underlying terminology used in this thesis to describe measurement quality, including the distinction between accuracy, precision, repeatability, and reproducibility \parencite{jcgm2012_international_vocabulary_of_metrology}.

Timing-measurement pitfalls have received particular attention. \textcite{chen2016_robust_benchmarking_noisy_environments} showed that consecutive run-time measurements are not necessarily independent and identically distributed, and recommended the minimum estimator over repeated short runs as a robust summary. \textcite{hoefler2015_scientific_benchmarking_parallel} reviewed 120 papers from three top conferences in high-performance computing and codified twelve rules for reporting parallel-system performance, emphasizing that distributions, not point averages, should be reported. \textcite{raasveldt2018_fair_benchmarking_database_systems} cataloged pitfalls specific to database benchmarking, including cold-versus-hot caches, untuned baselines, comparison of systems with different functionality, and ignored preprocessing and index-creation time, each of which can flip a published comparison.

Cloud environments add a further source of variation. \textcite{iosup2011_performance_variability_cloud_services} analyzed ten production services on \acrlong{aws} and Google App Engine over a year and found yearly, monthly, weekly, and daily patterns in service performance, with coefficients of variation above $1.10$ indicating high variability for several services. \textcite{schad2010_runtime_measurements_cloud} demonstrated that nominally identical \acrfull{ec2} instances differ measurably in \acrshort{cpu}, \acrshort{io}, and network throughput, with performance falling into two distinct bands tied to the underlying virtual hardware. \textcite{leitner2016_patterns_chaos_iaas} extended these results across four \acrshort{iaas} providers, including \acrshort{aws} \acrshort{ec2}, Google Compute Engine, Microsoft Azure, and IBM SoftLayer, codifying fifteen hypotheses about cloud performance variability and reporting that \acrshort{ec2} and Azure were substantially less predictable than the other two providers at the time of their study. \textcite{laaber2019_software_microbenchmarking_cloud} measured $4.5$ million microbenchmark runs across three public-cloud providers and found that running test and control experiments on the same instance in randomized order, combined with the Wilcoxon rank-sum test and bootstrapped confidence intervals, reliably detects slowdowns of $10\%$ or less.

Statistical comparison of systems is rarely reported with effect sizes. \textcite{arcuri2014_hitchhiker_statistical_tests} recommend reporting the Vargha--Delaney $\hat{A}_{12}$ effect size alongside non-parametric significance tests, with values near $0.5$ interpreted as no effect, $\geq 0.56$ as a small effect, $\geq 0.64$ as a medium effect, and $\geq 0.71$ as a large effect. The methodology chapter follows this approach.

